\begin{textsample}{NEG Dim 1 – 2019\_c4 – Score -5.00 – 2019\_c4/enade\_2019\_nutri\_acdc34d1d9.txt}  \label{ex:f1_neg_087}
CAPÍTULO 4 PERCEPÇÃO DA PROVA Nas análises feitas neste capítulo, constam as percepções dos concluintes da Área de Nutrição sobre a prova aplicada no Enade/2019.
Essas percepções foram mensuradas por meio de nove questões que avaliaram desde o grau de dificuldade da prova até o tempo gasto para concluí -la.
As percepções sobre a prova foram relacionadas ao desempenho dos estudantes, à Grande Região de funcionamento do curso, à Categoria Administrativa e ao tipo de Organização Acadêmica da IES.
O questionário de percepção da prova encontra -se ao final do Anexo VII, que traz a reprodução do exame.
O desempenho dos estudantes foi classificado em quatro quartos.
Para tanto, esse desempenho foi ordenado de forma ascendente.
O percentil 25, P25, também conhecido como primeiro quartil, é a nota de desempenho que deixa um quarto ( 25 \% ) dos valores observados abaixo e três quartos acima.
A Figura 1 apresenta uma ilustração desse conceito.
O quarto inferior de desempenho é composto pelas notas abaixo do primeiro quartil.
Já o percentil 75, P75, também conhecido como terceiro quartil, é o valor para o qual há três quartos ( 75 \% ) dos dados abaixo e um quarto acima dele.
O quarto superior de desempenho é composto pelas notas iguais ou acima do terceiro quartil.
O percentil 50, P50, também conhecido como mediana, é o valor que divide as notas em dois conjuntos de igual tamanho.
O segundo quarto inclui valores entre o primeiro quartil ( P25 ) e a mediana.
O terceiro quarto contém os valores entre a mediana ( P50 ) e o terceiro quartil ( P75 ).
Vale ressaltar que percentis, quartis e medianas são pontos que, não obrigatoriamente, pertencem ao conjunto original de dados, ao passo que os quartos são subconjuntos dos dados originais.
A seguir, serão apresentados gráficos com resultados selecionados, relativos às nove questões avaliadas por grupos de estudantes.
Nas barras dos gráficos, apresenta -se o percentual de alunos que assinalaram uma das opções ou a soma das porcentagens daqueles que assinalaram duas ( ou três ) delas.
Por exemplo, para as questões 1 e 2, nos gráficos, é apresentada a porcentagem total de participantes que assinalaram as opções ( D ) Difícil e ( E ) Muito difícil.
Em \textbf{cada} barra, foram assinalados, também, os extremos do intervalo de confiança de 95 \% como linhas verticais unidas por uma linha horizontal na forma da letra H maiúscula.
O estimador de um parâmetro com um \textbf{certo} nível de confiança ( e.g. 95,0 \% ) deve conter o parâmetro no intervalo de confiança em 95 \% das vezes.
Na comparação entre os estimadores dos parâmetros de duas classes de uma dada categoria ( e.g.
Norte e Nordeste nas Grandes Regiões, de primeiro e último quarto dentro de desempenho, de IES Públicas e Privadas, ou de Universidades e Faculdades ) associados aos seus respectivos intervalos de confiança, diz -se que não há diferença estatisticamente significativa entre os parâmetros das duas categorias quando há interseção entre os intervalos de confiança e que há diferença, se os intervalos de confiança são disjuntos ( para maiores informações vide Glossário ).
Nas Tabelas no Anexo II, são apresentados os valores absolutos e a distribuição percentual18 das alternativas válidas das nove questões, segundo o mesmo recorte de desempenho dos alunos e Grande Região de funcionamento do curso, de Categoria Administrativa e do tipo de Organização Acadêmica da IES. 4.1 GRAU DE DIFICULDADE DA PROVA 4.1.1 Componente de Formação Geral Ao avaliarem Qual o grau de dificuldade desta prova na parte de Formação Geral? ( Questão 1 ), 32,5 \% do grupo de inscritos e presentes optaram pelas alternativas Difícil ou Muito difícil.
Entretanto, para mais três quintos dos estudantes ( 61,7 \% ), o Componente de Formação Geral da prova foi considerado com grau de dificuldade Médio ( Gráfico 4.1, Gráfico 4.2 e, no Anexo II, a Tabela II.1 ).
O percentual de estudantes que consideraram a prova como Difícil ou Muito difícil foi maior na região Sul, onde a proporção foi de 37,2 \%, enquanto a de menor incidência foi a Nordeste, com 28,5 \%.
No Gráfico 4.1, é possível observar que as diferenças entre a região Norte e as regiões Centro-Oeste e Sul são estatisticamente significativas, assim como as diferenças entre a região Nordeste e as regiões Sudeste, Sul e Centro-Oeste, e, também, a diferença entre a região Sudeste e a região Sul.
Nas Grandes Regiões, a proporção de presentes à prova que consideraram o Componente de Formação Geral como sendo de grau de dificuldade Médio \textbf{esteve} entre 57,8 \%, na região Centro-Oeste e 65,4 \%, na região Nordeste. 18 Cumpre \textbf{lembrar} uma das convenções para tabelas numéricas ( pág. iii ) sobre a possibilidade de a soma das partes não resultar em 100 \% por questões de arredondamento. de dificuldade Médio \textbf{esteve} entre 57,8 \%, na região Centro-Oeste e 65,4 \%, na região Nordeste.
O percentual de alunos que consideraram o componente de Formação Geral da prova Difícil ou Muito difícil foi menor no quarto de melhor desempenho e maior no segundo e terceiro quartos : 31,5 \%, no primeiro quarto e 33,1 \%, no segundo e terceiro quartos, grupos de desempenho intermediário na prova.
No quarto de desempenho inferior, a proporção de alunos que consideraram a prova Difícil ou Muito difícil foi de 32,2 \%.
A diferença entre os quartos de desempenho não é estatisticamente significativa.
Para todos os quartos de desempenho, a alternativa modal para essa pergunta foi a Médio, com 62,8 \% e 61,0 \% dos respondentes nos quartos extremos, primeiro e quarto, respectivamente.
O percentual de estudantes que consideraram o componente de Formação Geral da prova como Difícil ou Muito difícil foi maior para os estudantes de Faculdades, para os quais a proporção foi de 34,7 \%, enquanto os CEFET/IF foram a Organização Acadêmica com a menor incidência, com 26,4 \%.
No Gráfico 4.3, é possível observar que há diferença estatisticamente significativa apenas entre Centros Universitários ( 31,7 \% ) e Faculdades.
Nas Universidades, a proporção de estudantes que consideraram o componente de Formação Geral da prova como Difícil ou Muito difícil foi de 32,2 \%.
Nas Organizações Acadêmicas, a proporção de presentes à prova que consideraram o Componente de Formação Geral como sendo de grau de dificuldade Médio \textbf{esteve} entre 60,9 \%, nas Faculdades e 67,9 \%, nos CEFET/IF ( ver também a Tabela II.2, no Anexo II ).
O percentual de estudantes que consideraram o componente de Formação Geral da prova como Difícil ou Muito difícil foi maior para os estudantes de cursos de IES Privadas ( 33,1 \% ), uma diferença estatisticamente significativa da proporção em IES Públicas ( 28,9 \% ).
Já a proporção de presentes à prova que consideraram o Componente de Formação Geral como sendo de grau de dificuldade Médio foi 61,8 \% para os de IES Públicas e 61,7 \% para os de IES Privadas ( ver Gráfico 4.4 ). 4.1.2 Componente de Conhecimento Específico Ao responderem à Questão 2 – Qual o grau de dificuldade desta prova na parte de Componente Específico? – 19,8 \% do grupo de estudantes classificaram -na como Difícil ou Muito difícil.
Além disso, o Componente modal de Conhecimento Específico da prova foi considerado com grau de dificuldade Médio por 71,1 \% dos alunos ( Gráfico 4.5, Gráfico 4.6, e, no Anexo II, a Tabela II.3 ).
A análise das respostas dos estudantes quanto ao grau de dificuldade do Componente de Conhecimento Específico da prova, agregado por Grande Região, mostra que a diferença entre a menor e as duas maiores proporções de alunos que a avaliaram como Difícil ou Muito difícil é estatisticamente significativa, sendo a menor, na região Nordeste ( 18,1 \% ), e as duas maiores, nas regiões Sul ( 22,3 \% ) e Centro-Oeste ( 21,9 \% ).
A diferença entre a região Nordeste e a região Sudeste ( 20,3 \% ) também é estatisticamente significativa.
O percentual de alunos que classificaram o grau de dificuldade como Médio, no Componente de Conhecimento Específico, variou de 68,4 \% a 73,0 \%, para as regiões Sul e Norte, respectivamente.
Considerando -se a avaliação da dificuldade das questões do Componente de Conhecimento Específico da prova, de acordo com o desempenho dos estudantes, observa -se que há diferença estatisticamente significativa dos resultados entre o primeiro quarto e os demais quartos de desempenho.
A proporção dos que classificaram a parte específica como Difícil ou Muito difícil variou de 16,9 \% ( último quarto ) a 23,9 \% ( primeiro quarto ).
As demais proporções dessa reposta foram de 19,8 \% e 18,9 \%, nos segundo e terceiro quartos, respectivamente.
Já a proporção dos que responderam que o grau de dificuldade das questões do Componente de Conhecimento Específico da prova foi Médio variou de 68,8 \%, no primeiro quarto a 72,6 \% no segundo.
Considerando -se a avaliação da dificuldade das questões do Componente de Conhecimento Específico da prova, o percentual de estudantes que as consideraram como Difícil ou Muito difícil foi maior para os estudantes de Faculdades, para os quais a proporção foi de 23,2 \%, enquanto os CEFET/IF foram a Organização Acadêmica com a menor incidência, com 9,4 \%.
No Gráfico 4.7, é possível observar que as diferenças entre as duas menores proporções, em CEFET/IF e Universidades ( 17,5 \% ), e as duas maiores proporções, em Faculdades e Centros Universitários ( 20,9 \% ), são estatisticamente significativas.
Nas Organizações Acadêmicas, a proporção de presentes à prova que consideraram o Componente de Formação Geral como sendo de grau de dificuldade Médio \textbf{esteve} entre 69,8 \%, nas Faculdades e 81,1 \%, nos CEFET/IF ( ver também Tabela II.4, no Anexo II ).
O percentual de estudantes que consideraram o Componente Específico da prova como Difícil ou Muito difícil foi maior para os estudantes de cursos de IES Privadas ( 21,2 \% ), uma diferença estatisticamente significativa da proporção em IES Públicas ( 12,7 \% ).
Já a proporção de presentes à prova que consideraram este componente como sendo de grau de dificuldade Médio foi 71,6 \% para os de IES Públicas e 71,0 \% para os de IES Privadas ( ver Gráfico 4.8 e Tabela II.4, no Anexo II ). 4.2 EXTENSÃO DA PROVA EM RELAÇÃO AO TEMPO TOTAL Indagados quanto à extensão da prova, em relação ao tempo total oferecido para a sua resolução ( Questão 3 ), os estudantes apontaram, com maior incidência, a alternativa que considerava a extensão Adequada, para todas as agregações consideradas ( Gráfico 4.9, Gráfico 4.10, e, no Anexo II, a Tabela II.5 ).
O percentual de alunos que responderam ser a extensão da prova Adequada foi de 63,2 \%.
Já 31,4 \% dos inscritos presentes consideraram que a prova foi Longa ou Muito longa, e 5,4 \% a avaliaram como Curta ou Muito curta.
Dentre as Grandes Regiões, a proporção daqueles que avaliaram a prova como Longa ou Muito longa, em relação ao tempo total destinado à sua resolução, variou de 27,7 \% na região Norte até 32,5 \%, na região Sul.
Há diferença estatisticamente significativa entre as regiões Nordeste e Centro-Oeste, e também entre a região Norte e as regiões Nordeste, Sudeste e Sul.
Considerando -se o desempenho dos alunos, nota -se, ainda, que 60,4 \% consideraram a extensão da prova Adequada, no quarto de desempenho inferior e 65,5 \%, no de melhor desempenho ( quarto superior ).
Nos quartos intermediários, essa proporção foi de 62,9 \%, no segundo e 63,9 \%, no terceiro quarto, proporções crescentes com o desempenho.
No Gráfico 4.6, pode -se constatar que há diferença estatisticamente significativa entre as proporções de estudantes que consideraram a prova Longa ou Muito longa em função dos quartos de desempenho, do primeiro quarto em relação ao terceiro e último quartos, e do segundo em relação ao quarto de melhor desempenho.
As proporções são decrescentes em função do desempenho, sendo a maior proporção dessa resposta no primeiro quarto de desempenho ( 34,4 \% ), e a menor, no quarto superior ( 28,9 \% ).
Considerando -se a extensão da prova, o percentual de estudantes que avaliaram a prova como Longa ou Muito longa, em relação ao tempo total destinado à sua resolução, foi maior para os estudantes de Centros Universitários, para os quais a proporção foi de 34,0 \%, enquanto os CEFET/IF foram a Organização Acadêmica com a menor incidência, com 15,1 \%.
No Gráfico 4.11, é possível observar que as diferenças entre todas as organizações são estatisticamente significativas.
Nas Organizações Acadêmicas, a proporção de presentes à prova que consideraram a sua extensão como Adequada, em relação ao tempo total destinado à sua resolução, \textbf{esteve} entre 60,6 \%, nos Centros Universitários e 83,0 \% nos CEFET/IF ( ver também Tabela II.6, no Anexo II ).
O percentual de estudantes que consideraram a prova Longa ou Muito longa, em relação ao tempo total destinado à sua resolução, foi maior para os estudantes de cursos de IES Privadas ( 32,2 \% ), uma diferença estatisticamente significativa da proporção em IES Públicas ( 26,6 \% ).
Já a proporção de presentes à prova que consideraram a extensão como Adequada foi 69,0 \% para os de IES Públicas e 62,1 \% para os de IES Privadas ( ver Gráfico 4.12 e Tabela II.6, no Anexo II ). 4.3 COMPREENSÃO DOS ENUNCIADOS DAS QUESTÕES 4.3.1 Componente de Formação Geral Em relação aos enunciados das questões do Componente de Formação Geral ( Questão 4 ), as opiniões foram positivas, já que 72,0 \% dos alunos avaliados consideraram todos ou a maioria dos enunciados das questões claros e objetivos ( Gráfico 4.13, Gráfico 4.14, e, no Anexo II, a Tabela II.7 ).
Na análise regional, a percentagem de estudantes que avaliaram que todos ou a maioria dos enunciados das questões do Componente de Formação Geral \textbf{estavam} claros e objetivos variou de 68,7 \%, na região Norte a 75,2 \%, na região Sul.
A análise das percepções dos estudantes sobre a clareza e a objetividade dos enunciados permite afirmar que todos, ou a maioria dos enunciados de questões relativas ao Componente de Formação Geral, foram considerados claros e objetivos para a maior parte dos respondentes ( maior ou igual a 68,7 \%, em todas as regiões e maior ou igual a 67,8 \% para todos os quartos de desempenho ).
Segundo o desempenho, observa -se que a proporção dos que emitiram essa opinião cresce conforme o desempenho aumenta, com diferenças estatisticamente significativas entre todos os quartos de desempenho, com exceção da existente entre os dois primeiros quartos.
No quarto superior, a clareza e a objetividade de todos ou da maioria dos enunciados das questões foram percebidas por 76,6 \%.
A proporção dos presentes que consideraram que todos ou a maioria dos enunciados das questões do Componente de Formação Geral \textbf{estavam} claros e objetivos foi maior para os estudantes de Universidades, para os quais a proporção foi de 72,7 \%, enquanto as Faculdades foram a Organização Acadêmica com a menor incidência, com 69,4 \%.
No Gráfico 4.15, é possível observar que as diferenças entre Faculdades, a menor proporção, e Universidades e Centros Universitários ( 72,5 \% ), as duas maiores proporções, são estatisticamente significativas ( ver também Tabela II.8, no Anexo II ).
O percentual de estudantes que consideraram que todos ou a maioria dos enunciados das questões do Componente de Formação Geral \textbf{estavam} claros e objetivos foi maior para os estudantes de cursos de IES Públicas ( 73,7 \% ), sem diferença estatisticamente significativa da proporção em IES Privadas ( 71,7 \% ).
Um maior \textbf{detalhamento} pode ser encontrado no Gráfico 4.16 e na Tabela II.8, no Anexo II. 4.3.2 Componente de Conhecimento Específico Em relação aos enunciados das questões do Componente de Conhecimento Específico da prova, para 85,7 \% dos estudantes avaliados da Área de Nutrição, a clareza e a objetividade ( Questão 5 ) \textbf{estavam} presentes em todas ou na maioria das questões ( Gráfico 4.17 a Gráfico 4.20, e no Anexo II, as Tabelas II.9 e II.10 ).
A maioria dos estudantes de todas as Grandes Regiões brasileiras considerou claros e objetivos todos ou a maioria dos enunciados das questões do Componente de Conhecimento Específico da prova, percentual sempre maior ou igual a 81,4 \%.
As diferenças entre a região Norte e as regiões Sul e Centro-Oeste são estatisticamente significativas, bem como entre a região Nordeste e as regiões Sudeste e Sul.
A proporção de estudantes que consideraram os enunciados das questões claros e objetivos apresenta uma tendência crescente em relação ao aumento de desempenho : mais elevada no quarto superior ( 92,7 \% ), se comparada ao quarto inferior de desempenho ( 76,6 \% ).
As diferenças entre todos os quartos de desempenho são estatisticamente significativas.
A proporção dos presentes que consideraram que todos ou a maioria dos enunciados das questões do Componente de Conhecimento Específico \textbf{estavam} claros e objetivos foi maior para os estudantes de CEFET/IF, para os quais a proporção foi de 90,6 \%, enquanto as Faculdades foram a Organização Acadêmica com a menor incidência, com 82,5 \%.
No Gráfico 4.19, é possível observar que as diferenças entre Faculdades, Centros Universitários ( 84,9 \% ) e Universidades ( 87,7 \% ) são estatisticamente significativas ( ver também Tabela II.10, no Anexo II ).
O percentual de estudantes que consideraram que todos ou a maioria dos enunciados das questões do Componente de Conhecimento Específico \textbf{estavam} claros e objetivos foi maior para os estudantes de cursos de IES Públicas ( 90,6 \% ), uma diferença estatisticamente significativa da proporção em IES Privadas ( 84,8 \% ).
Ver Gráfico 4.20 e Tabela II.10, no Anexo II, para um maior \textbf{detalhamento}. 4.4 SUFICIÊNCIA DAS INFORMAÇÕES/INSTRUÇÕES FORNECIDAS Ao avaliarem as informações/instruções fornecidas para a resolução das questões ( Questão 6 ), 88,5 \% dos respondentes da Área de Nutrição de todo o Brasil afirmaram que estas eram até excessivas ou suficientes em todas ou na maioria das questões ( Gráfico 4.21, Gráfico 4.22, e, no Anexo II, a Tabela II.11 ).
Quanto à distribuição de respondentes pelas Grandes Regiões, observa -se que a proporção de estudantes que consideraram as informações/instruções fornecidas até excessivas ou suficientes em todas ou na maioria das questões foi sempre superior ou igual a 83,7 \%, chegando a 90,6 \%, na região Sul.
Existem diferenças estatisticamente significativas entre as regiões Norte e Nordeste e as demais regiões.
Levando -se em conta o desempenho dos participantes, observa -se uma tendência crescente das proporções de participantes que avaliaram as informações/instruções como até excessivas ou suficientes em todas ou na maioria das questões, havendo diferenças estatisticamente significativas entre todos os quartos de desempenho, como mostrado no Gráfico 4.12.
O percentual foi mais elevado no quarto superior ( 92,9 \% ), percentual superior à média nacional ( 88,5 \% ).
Já no quarto inferior, a suficiência das informações/instruções declarada como até excessiva, em todas ou na maioria das questões, foi percebida por 82,6 \% dos respondentes.
Considerando -se as informações/instruções fornecidas para a resolução das questões ( Questão 6 ), a proporção que considerou que estas eram até excessivas ou suficientes em todas ou na maioria das questões foi maior para os estudantes de CEFET/IF, para os quais a proporção foi de 94,3 \%, enquanto as Faculdades foram a Organização Acadêmica com a menor incidência, com 85,1 \%.
No Gráfico 4.23, é possível observar que as diferenças entre Faculdades e as demais organizações acadêmicas são estatisticamente significativas ( ver também Tabela II.12, no Anexo II ).
O percentual de estudantes que consideraram que as informações/instruções fornecidas para a resolução das questões eram até excessivas ou suficientes em todas ou na maioria das questões foi maior para os estudantes de cursos de IES Públicas ( 92,6 \% ), uma diferença estatisticamente significativa da proporção em IES Privadas ( 87,7 \% ) ( ver Gráfico 4.24 e Tabela II.12, no Anexo II ). 4.5 DIFICULDADE ENCONTRADA AO RESPONDER À PROVA Indagados sobre as dificuldades com as quais se depararam ao responder à prova ( Questão 7 ), 16,2 \% dos estudantes apontaram o Desconhecimento do conteúdo.
Para 46,3 \%, a Forma diferente de abordagem do conteúdo foi indicada como dificuldade.
Já a Falta de motivação para fazer a prova foi a dificuldade apontada por 9,3 \% dos respondentes.
Considerando -se todo o Brasil, 24,6 \% dos respondentes afirmaram que não tiveram qualquer tipo de dificuldade para responder à prova ( Tabelas II.13 e II.14, no Anexo II ).
Nos Gráficos 4.25 a 4.28, são apresentados os percentuais de estudantes que apontaram o Desconhecimento do conteúdo como dificuldade ao responder à prova.
Na análise por Grandes Regiões, o percentual de inscritos e presentes que apontaram o Desconhecimento do conteúdo como dificuldade ao responder à prova não superou 18,4 \%.
Os percentuais variaram de 15,0 \%, na região Centro-Oeste a 18,4 \%, na região Sul.
A Forma diferente de abordagem do conteúdo foi a escolha modal dos estudantes, com percentuais que variaram de 44,4 \% ( região Nordeste ) a 47,8 \% ( região Centro-Oeste ).
O percentual de alunos que citaram a Falta de motivação como dificuldade variou de 8,5 \% ( região Sudeste ) a 10,3 \% ( região Norte ).
Os que declararam não ter qualquer tipo de dificuldade para responder à prova variaram de 21,5 \%, na região Norte a 26,4 \%, na região Nordeste.
Em relação aos quartos de desempenho, o Desconhecimento do conteúdo foi a opção escolhida por 15,3 \% dos estudantes do quarto inferior e por 16,7 \% do quarto superior, caracterizando uma tendência crescente.
Não são observadas diferenças estatisticamente significativas entre os quartos de desempenho.
A alternativa modal para os alunos, quando agregados pelos quartos de desempenho, para a dificuldade encontrada, foi causada pela Forma diferente de abordagem do conteúdo : 48,5 \%, no quarto inferior e 43,0 \%, do último quarto, escolheram essa alternativa.
Na análise por Organização Acadêmica, o percentual de inscritos e presentes que apontaram o Desconhecimento do conteúdo como dificuldade ao responder à prova não superou 18,9 \%.
Os percentuais variaram de 15,4 \%, para Faculdades a 18,9 \%, nos CEFET/IF.
A Forma diferente de abordagem do conteúdo foi a escolha modal dos estudantes, com percentuais que variaram de 44,1 \% ( Universidades ) a 49,9 \% ( Faculdades ).
O percentual de alunos que citaram a Falta de motivação como dificuldade variou de 7,3 \% ( Faculdades ) a 10,7 \% ( Universidades ).
Os que declararam não ter qualquer tipo de dificuldade para responder à prova variaram de 22,6 \%, nos CEFET/IF a 25,0 \%, nos Centros Universitários.
No Gráfico 4.27, é possível observar que as diferenças entre as organizações acadêmicas não são estatisticamente significativas ( ver também Tabela II.14, no Anexo II ).
Em relação às Categorias Administrativas, o Desconhecimento do conteúdo foi a opção escolhida por 16,8 \% dos estudantes nas IES Públicas e por 16,1 \%, nas IES Privadas.
Não são observadas diferenças estatisticamente significativas entre estas duas categorias.
A alternativa modal para os alunos, quando agregados pelas Categorias Administrativas, para a dificuldade encontrada foi causada pela Forma diferente de abordagem do conteúdo : 39,0 \%, para as IES Públicas e 47,7 \%, para as IES Privadas ( ver Gráfico 4.28 e Tabela II.14, no Anexo II ). 4.6 CONTEÚDOS DAS QUESTÕES OBJETIVAS DA PROVA Ao analisar os conteúdos das questões objetivas da prova ( Questão 8 ), um percentual muito pequeno dos estudantes avaliados, apenas 2,0 \%, afirmaram que não estudou ainda a maioria desses conteúdos ( Gráficos 4.29 a Gráfico 4.32, e as Tabelas II.15 e II.16, no Anexo II ).
A maioria ( 80,9 \% ) afirmou ter estudado e aprendido muitos ou todos os conteúdos avaliados.
Na análise por Grande Região, a proporção de respondentes que escolheram a opção não estudou ainda a maioria desses conteúdos foi pequena.
Observa -se que, nas regiões Norte ( 3,8 \% ) e Nordeste ( 2,5 \% ), as proporções foram maiores do que a média nacional ( 2,0 \% ).
Nas regiões Sudeste ( 1,6 \% ), Sul ( 1,3 \% ) e Centro-Oeste ( 1,5 \% ), a situação foi inversa : uma proporção mais baixa.
São observadas diferenças estatisticamente significativas das regiões Norte e Nordeste, em relação às regiões Sudeste, Sul e Centro-Oeste.
Em todas as regiões, a maioria dos presentes afirmou ter estudado e aprendido muitos ou todos os conteúdos, com proporções variando entre 70,1 \%, na região Norte a 84,4 \%, na região Centro-Oeste.
Considerando -se separadamente as opiniões de estudantes dos quatro quartos de desempenho, observa -se que, no quarto inferior, 4,4 \% ofereceram como resposta que não estudou ainda a maioria desses conteúdos, sendo 0,5 \% os do quarto superior com a mesma resposta.
A diferença entre os alunos que optaram por esse motivo de dificuldade nos quartos inferiores, em relação aos quartos superiores de desempenho, é estatisticamente significativa.
Tendo -se em conta o quarto superior, 91,4 \% dos alunos afirmaram ter estudado e aprendido muitos ou todos os conteúdos.
No outro extremo, no primeiro quarto, 68,8 \% optaram pelas mesmas categorias.
As proporções são crescentes com o desempenho.
Na análise por Organização Acadêmica, a proporção de respondentes que escolheram a opção não estudou ainda a maioria desses conteúdos foi pequena.
Observa -se que, nas Faculdades ( 2,7 \% ) e Centros Universitários ( 2,3 \% ), as proporções foram maiores do que a média nacional ( 2,0 \% ).
Nos CEFET/IF ( 1,9 \% ) e Universidades ( 1,5 \% ), a situação foi inversa : uma proporção mais baixa.
No Gráfico 4.31, é possível observar que as diferenças de Universidades para Faculdades e Centros Universitários são estatisticamente significativas ( ver também Tabela II.16, no Anexo II ).
Em todos os tipos de Organização Acadêmica, a maioria dos presentes afirmou ter estudado e aprendido muitos ou todos os conteúdos, com proporções variando entre 76,1 \%, nas Faculdades a 83,7 \%, nas Universidades.
Em relação às Categorias Administrativas, observa -se que, nas IES Públicas, 1,0 \% ofereceu como resposta que não estudou ainda a maioria desses conteúdos, sendo 2,2 \% nas IES Privadas com a mesma resposta.
A diferença entre os alunos que optaram por esse motivo de dificuldade nas duas categorias é estatisticamente significativa.
Considerando -se as IES Públicas, 87,6 \% dos alunos afirmaram ter estudado e aprendido muitos ou todos os conteúdos.
Na outra categoria, 79,7 \% optaram pelas mesmas categorias ( ver Gráfico 4.32 e Tabela II.16, no Anexo II ). 4.7 TEMPO GASTO PARA CONCLUIR A PROVA Ao responderem sobre o tempo de conclusão da prova ( Questão 9 ), mais de quatro quintos dos estudantes ( 83,3 \% ) afirmaram ter gasto Entre duas e quatro horas ( Gráfico 4.33 a Gráfico 4.36 e, no Anexo II, as Tabelas II.17 e II.18 ).
Considerando -se as cinco Grandes Regiões brasileiras, os percentuais dos que utilizaram Entre duas e quatro horas para finalizar a prova nas regiões Norte ( 84,9 \% ) e Nordeste ( 84,6 \% ) foram superiores ao percentual nacional.
Nas demais Grandes Regiões, o percentual de alunos que dispensaram Entre duas e quatro horas para concluir a prova foi igual ou menor que a média : 83,3 \%, na região Sudeste, 82,3 \%, na região Centro-Oeste e 79,5 \%, na região Sul, como mostrado no Gráfico 4.33.
A diferença entre a região Sul e as regiões Norte, Nordeste e Sudeste é estatisticamente significativa.
Uma vez analisadas as alternativas escolhidas pelos estudantes que se situam nos diferentes quartos de desempenho, observa -se uma tendência crescente : uma proporção maior de participantes no quarto superior declarou ter gasto Entre duas e quatro horas para concluir a prova quando comparada com a proporção de participantes nos quartos inferiores.
Não há diferença estatisticamente significativa apenas entre o terceiro e o último quartos de desempenho.
Na análise por Organização Acadêmica, o percentual dos que utilizaram Entre duas e quatro horas para finalizar a prova nas Universidades ( 82,5 \% ) foi inferior ao percentual nacional.
Nos demais tipos de Organização Acadêmica, o percentual de alunos que dispensaram Entre duas e quatro horas para concluir a prova foi maior do que a média : 83,5 \%, nas Faculdades, 84,2 \%, nos Centros Universitários e 84,9 \%, nos CEFET/IF, como mostrado no Gráfico 4.35.
A diferença entre Centros Universitários e Universidades é estatisticamente significativa.
Em relação às Categorias Administrativas, observa -se que, nas IES Públicas, 80,0 \% declararam ter gasto Entre duas e quatro horas para concluir a prova, sendo 83,9 \%, nas IES Privadas, com a mesma resposta.
A diferença entre os alunos que optaram por esse motivo de dificuldade nas duas categorias é estatisticamente significativa ( ver Gráfico 4.36 e Tabela II.18, no Anexo II ).

% matched lemmas: cada, certo, detalhamento, estar, lembrar
\end{textsample}
